% Figure from MATH 452 notes, section 22. Compile with the notes preamble.
\begin{tikzpicture}[scale=1.0]
% left: a star-shaped domain
\filldraw[fill=figB!10, draw=figB, thick] plot[smooth cycle, tension=0.55] coordinates {(0.000,1.350) (-0.411,0.566) (-1.284,0.417) (-0.666,-0.216) (-0.794,-1.092) (-0.000,-0.700) (0.794,-1.092) (0.666,-0.216) (1.284,0.417) (0.411,0.566)};
\coordinate (p0) at (0.05,-0.05);
\draw[figB!70, thin] (p0) -- (0.000,1.256);
\draw[figB!70, thin] (p0) -- (-0.370,0.510);
\draw[figB!70, thin] (p0) -- (-1.194,0.388);
\draw[figB!70, thin] (p0) -- (-0.599,-0.195);
\draw[figB!70, thin] (p0) -- (-0.738,-1.016);
\draw[figB!70, thin] (p0) -- (-0.000,-0.630);
\draw[figB!70, thin] (p0) -- (0.738,-1.016);
\draw[figB!70, thin] (p0) -- (0.599,-0.195);
\draw[figB!70, thin] (p0) -- (1.194,0.388);
\draw[figB!70, thin] (p0) -- (0.370,0.510);
\filldraw[black] (p0) circle (1.4pt) node[anchor=north west, inner sep=1.5pt, font=\scriptsize] {$\mathbf{p}_0$};
\node[font=\scriptsize, align=center] at (0,-1.95) {star-shaped: all segments\\ from $\mathbf{p}_0$ stay inside\\ closed $\Rightarrow$ exact};
% right: the punctured plane (a disk of it shown)
\begin{scope}[shift={(4.3,0)}]
\fill[figB!10] (0,0) circle (1.35);
\draw[black!35, thin, dashed] (0,0) circle (1.35);
\coordinate (p0) at (-0.75,0.3);
\coordinate (q) at (0.9,-0.36);
\draw[figB!70, thin] (p0) -- (0.3,1.05);
\draw[figB!70, thin] (p0) -- (-1.15,-0.55);
\draw[figR, thick] (p0) -- (q);
\filldraw[fill=white, draw=figR, thick] (0,0) circle (2.2pt);
\node[figR, font=\scriptsize, anchor=north, inner sep=1pt] at (-0.3,-0.12) {$\mathbf{0}$ removed};
\filldraw[black] (p0) circle (1.4pt) node[anchor=south east, inner sep=1.5pt, font=\scriptsize] {$\mathbf{p}_0$};
\filldraw[black] (q) circle (1.2pt) node[anchor=west, inner sep=1.5pt, font=\scriptsize] {$\mathbf{p}$};
\node[font=\scriptsize, align=center] at (0,-1.95) {$\mathbb{R}^2\setminus\{\mathbf{0}\}$: the segment\\ from $\mathbf{p}_0$ to $\mathbf{p}$ hits the hole\\ closed $\not\Rightarrow$ exact};
\end{scope}
\end{tikzpicture}%
